\scanpage{1}
% not manuscript: couverture de l'Université de Montpellier / IMAG, « La "Longue Marche" à travers la théorie de Galois », cote n° 140-1, 36 pages.

\scanpage{2}
% S002-E01: upper right heading.
\margin{\uline{``Long March''}}
% S002-E02: upper left, cancelled beginning and equation.
\[
\deleted{\underline{u}\unsure{\Rightarrow|=}}\unsure{\,\cdot\,|\,;\,}
\qquad \alpha=\underline{u}a^{-1}
\]
% S002-E03: detached expression under the heading.
\[
-\unsure{v|\nu}d^2+wcd+c^2
\]
% S002-E04: upper left.
\[
\underline{u}=\alpha\underline{a}
\]
% S002-E05: upper right.
\[
\det\bigl(d\rho+\unsure{c}\bigr)
=d^2\det\rho+cd\,\uline{\operatorname{Tr}\rho}+c^2
\]
% S002-E06: upper left matrix, without a label in the original.
\[
\begin{pmatrix}a&b\\c&d\end{pmatrix}
\]
% S002-E07: upper right matrix.
\[
\rho=\begin{pmatrix}0&\unsure{v_1|\nu_1}\\\unsure{v|\nu}&w\end{pmatrix}
\]
% S002-E08: left.
\[
a=\alpha^{-1}u
\]
% S002-E09: right.
\[
d_0\rho+c_0=
\begin{pmatrix}
c_0&d_0\unsure{v_1|\nu_1}\\
d_0\unsure{v|\nu}&d_0w+c_0
\end{pmatrix}
\]
% S002-E10: middle. The final matrix is left unfinished in the scan.
\[
\begin{pmatrix}a&b\\c&d\end{pmatrix}
\begin{pmatrix}c&d\\d\unsure{v|\nu}&dw+c\end{pmatrix}
=\left(\begin{array}{cc}ac+bd\unsure{v|\nu}&\\&\end{array}\right.
\]
% S002-E11: middle, overwritten top entries in the second factor.
\[
\underline{a}=\alpha^{-1}u
=\frac{1}{\delta_\alpha}
\begin{pmatrix}d&-b\\-c&a\end{pmatrix}
\begin{pmatrix}
\deleted{\unsure{A|B}}\inserted{\unsure{\mu}}&
\deleted{\unsure{B}}\inserted{\unsure{\lambda}}\\
0&1
\end{pmatrix}
\]
% S002-E12: continuation across the fold.
\[
=\frac{1}{\delta_\alpha}
\begin{pmatrix}
\deleted{\unsure{d|\delta}}\inserted{d}\mu&d\lambda-b\\
-c\mu&-c\lambda+a
\end{pmatrix}
=d_0\rho+\unsure{c_0|d_0}
\]
% S002-E13: lower middle.
\[
c_0=d\mu\unsure{/|\ne}\delta_\alpha
=\frac{-d\unsure{v|\nu}}{-\unsure{v^2|\nu^2}+wcd+\unsure{c^2}}
\]
% S002-E14: cancelled line; no closing parenthesis is written.
\[
\deleted{
d_0=\frac{1}{\delta_\alpha}(d\lambda-b)
=\frac{1}{\delta_\alpha}
\left(-\unsure{v|\nu}bd^2+\unsure{v_1|\nu_1}acd+w bcd\right.
}
\]
% S002-E15: immediately below the cancellation.
\[
d_0\unsure{v|\nu}=\frac{-c\mu}{\delta_\alpha}
\]
% S002-E16: abandoned beginning at left.
\[
d_0=-
\]
% S002-E17: lower right.
\[
c_0=\frac{d\mu}{\delta_\alpha}
\]
% S002-E18: lower right.
\[
d_0=\frac{-c\mu}{\delta_\alpha\unsure{v|\nu}}
\]
% S002-E19: cancelled calculation to the left of the box.
\[
\deleted{\det\underline{a}=\det\unsure{\,\cdot\,}}
\qquad
\deleted{\frac{\unsure{u|\mu}}{\unsure{\delta_\alpha|\delta}}}
\]
% S002-E20: cancelled mark beside the box.
\margin{\(\deleted{\unsure{\underline{a}}}\)}
% S002-E21: boxed expression, including the ambiguous sign after the parenthesis.
\[
\boxed{\underline{a}=\frac{\mu}{\delta_\alpha}
\left(\unsure{-|=}\frac{c}{\unsure{v|\nu}}\rho+d\right)}
\]
% S002-E22: lower left; the question mark is written above the equals sign.
\[
\det a\overset{?}{=}\delta_\alpha^{-1}\mu
\]
% S002-E23: bottom right. The terminal dots are present in the scan.
\[
\det\underline{a}=\left(\frac{\mu}{\delta_\alpha}\right)^2
\left[\frac{c^2}{\unsure{v^2|\nu^2}}
\bigl(-\unsure{vv_1|\nu\nu_1}\bigr)
\unsure{-|+}\frac{dc}{\unsure{v|\nu}}w+d^2\right]\ldots
\]
% S002-E24: lower-left pencil foliation.
\margin{1}

\scanpage{3}
% not manuscript: circulaire dactylographiée de l'Université des Sciences et Techniques du Languedoc, 18 septembre 1978, « Rentrée universitaire 1978-1979 », adressée au professeur Grothendieck; feuille à l'envers dans le scan.

\scanpage{4}
\mspage{1}

\[
\unsure{\widetilde{\Pi}_n^{\circ D}|\widetilde{\Pi}_n^D}
\unsure{\simeq|=}
\deleted{\unsure{\{}}
\{\tilde\rho,\tilde\sigma\mid\tilde\rho^n=\tilde\sigma^2\}
\qquad
\boxed{\tilde\rho^n=\tilde\sigma^2=\tilde\omega_0}
\quad\text{central \unsure{element}}
\]

\diagram{\begin{tikzcd}[ampersand replacement=\&,column sep=small]
1\arrow[r]\&\mathbb Z\arrow[r]\&
\unsure{\widetilde{\Pi}_n^{\circ D}|\widetilde{\Pi}_n^D}\arrow[r,no head]\&
\Pi_n^D\arrow[r,no head]\&1
\end{tikzcd}}
central extension

Let us put
\[
\begin{aligned}
\boxed{\tilde\rho_0=\tilde\sigma\tilde\rho\tilde\sigma^{-1}\tilde\rho}
&=\tilde\sigma(\tilde\rho)\tilde\rho\\
&=(\tilde\sigma\tilde\rho)^2\tilde\omega_0^{-1}
\end{aligned}
\qquad
\boxed{\tilde\rho_i=\tilde\rho^{\unsure{i}}(\tilde\rho_0)}
\]
\margin{\deleted{$\unsure{+}$} $\tilde\rho_i$ lifts $\rho_i$}

We have
\[
\begin{aligned}
\tilde\rho_{n-1}\tilde\rho_{n-2}\cdots\tilde\rho_0
&=\underbrace{\tilde\rho^{n-1}\tilde\rho_0\tilde\rho^{-(n-1)}}
\underbrace{\tilde\rho^{n-2}\tilde\rho_0\tilde\rho^{-(n-2)}}
\cdots(\tilde\rho\tilde\rho_0\tilde\rho^{-1})\tilde\rho_0\\
&=\tilde\rho^{n-1}\tilde\rho_0\tilde\rho^{-1}\tilde\rho_0\tilde\rho^{-1}
\cdots\tilde\rho^{-1}\deleted{\unsure{{}_0}}\tilde\rho_0\tilde\rho^{-1}\tilde\rho_0
\qquad\begin{gathered}(n\ \text{\unsure{factors}}\\\tilde\rho_0\end{gathered}\\
&=\tilde\rho^n\underbrace{(\tilde\rho^{-1}\tilde\rho_0)\cdots(\tilde\rho^{-1}\tilde\rho_0)}_{n\ \text{factors}}
=\underbrace{\tilde\rho^n}_{\tilde\omega_0}(\tilde\rho^{-1}\tilde\rho_0)^n
\end{aligned}
\]

now
\[
\tilde\rho^{-1}\tilde\rho_0
=\tilde\rho^{-1}\tilde\sigma\tilde\rho\tilde\sigma^{-1}\tilde\rho
=(\tilde\rho^{-1}\tilde\sigma)(\tilde\rho)
\]
. \unsure{therefore}
\[
(\tilde\rho^{-1}\tilde\rho_0)^n
=(\tilde\rho^{-1}\tilde\sigma)(\tilde\rho^n)
=(\tilde\rho^{-1}\tilde\sigma)(\tilde\omega_0)
=\tilde\omega_0
\]
whence
\[
\boxed{\tilde\rho_{n-1}\tilde\rho_{n-2}\cdots\tilde\rho_0=\tilde\omega_0^2}
\]

Let $\widetilde\Pi_n$ be
\deleted{the \unsure{subgroup} of $\unsure{\widetilde\Pi_n^D|\Pi_n^D}$ \unsure{generated}}
\inserted{the inverse image of $\Pi_n\subset\Pi_n^D$ in $\widetilde\Pi_n^D$}
\deleted{by the $\tilde\rho_i$, it is \unsure{invariant} under $\rho$, and also under $\sigma$}
\margin{\deleted{$\unsure{\widetilde\Pi_n^D}\ \unsure{\neq}\ \tilde\rho_{n-1}\tilde\rho_{n-2}$}}

\diagram{\begin{tikzcd}[ampersand replacement=\&,column sep=small,row sep=small]
1\arrow[r,no head]\&\mathbb Z\arrow[r]\&\widetilde\Pi_n\arrow[r,no head]\arrow[d,equal]\&\Pi_n\arrow[r,no head]\&1\\
\&\&\left\{\tilde\rho_0,\ldots,\tilde\rho_{n-1},\tilde\omega_0\ \middle|\ 
\begin{array}{l}\tilde\omega_0\ \text{central},\\
\tilde\rho_{n-1}\tilde\rho_{n-2}\cdots\tilde\rho_0=\tilde\omega_0^2
\end{array}\right\}
\end{tikzcd}}

Let
\margin{$\unsure{\tilde\rho|\rho}$}
\[
\widetilde\Pi_n^\circ=
\left\{\tilde\rho_0,\ldots,\tilde\rho_{n-1}\ \middle|\ 
\begin{array}{l}
\tilde\rho_{n-1}\tilde\rho_{n-2}\cdots\tilde\rho_0\ \text{central}\quad\text{i.e.}\\
\tilde\rho_{n-1}\tilde\rho_{n-2}\cdots\tilde\rho_0
=\tilde\rho_{n-2}\tilde\rho_{n-3}\cdots\tilde\rho_0\tilde\rho_{n-1}=\cdots
\end{array}\right\}
\]

Then
\diagram{\begin{tikzcd}[ampersand replacement=\&,column sep=small]
1\arrow[r,no head]\&\mathbb Z\arrow[r,no head]\&\widetilde\Pi_n^\circ\arrow[r,no head]\&
\overset{\deleted{\unsure{\mathrm{can}}}}{\Pi_n}\arrow[r,no head]\&1
\end{tikzcd}}
\unsure{generated} by
\[
\bar\omega=\tilde\rho_{n-1}\tilde\rho_{n-2}\cdots\tilde\rho_0
\]
\[
\widetilde\Pi_n^{\deleted{\unsure{\circ}}}
\unsure{\simeq|\subset}
\widetilde\Pi_n^\circ\mathop{\wedge}_{\mathbb Z\bar\omega}\mathbb Z\omega
\]
\diagram{\begin{tikzcd}[ampersand replacement=\&,column sep=small]
\bar\omega_0\arrow[r]\&\omega_0^2
\end{tikzcd}}
\margin{3}

\scanpage{5}
% not manuscript: sujet d'examen imprimé, session de Juin 1978, C4 Math. Pures et Approfondies, option de M. Petrich; question sur les classes de conjugaison et l'icosaèdre.

\scanpage{6}
\mspage{2}

$\widetilde\Pi_n^\circ$ normal in $\widetilde\Pi_n$
(\unsure{subgroup} of index 2)
\margin{\unsure{subgroup} \unsure{generated} by the $\tilde\rho_i$}

is it normal in $\widetilde\Pi_n^D$?
\unsure{Invariant|Invariance} under $\tilde\rho$
\unsure{it is} clear --- hence normal in $\widetilde\Pi_n^{D+}$.
\unsure{Invariant|Invariance} under $\tilde\sigma$?

in $\Pi_n^D$
\[
\begin{aligned}
\sigma^{\unsure{\ast|-1}}(\rho_0)&=\rho_1,\quad
\sigma^{\unsure{\ast|-1}}(\rho_1)=\rho_0,\quad
\sigma^{\unsure{\ast|-1}}(\rho_2)=\rho_0(\rho_{n-1}),\quad\cdots\\
\sigma^{\unsure{\ast|-1}}(\rho_i)&=\rho_n\cdots\rho_{n+2-i}(\rho_{n+1-i})\cdots
\end{aligned}
\]
in $\widetilde\Pi_n^D$
\[
\begin{aligned}
\tilde\sigma^{\unsure{\ast|-1}}(\tilde\rho_0)&=\tilde\rho_1\tilde z_0,\quad
\tilde\sigma^{\unsure{\ast|-1}}(\tilde\rho_1)=\tilde\rho_0\tilde z_1\quad\cdots\\
\tilde\sigma^{\unsure{\ast|-1}}(\tilde\rho_i)&=
(\tilde\rho_n\cdots\tilde\rho_{\unsure{n+2-i}})(\tilde\rho_{n+1-i})\unsure{\tilde z_i|\hat z_i}
\end{aligned}
\]
\[
z_{n-1}z_{n-2}\cdots z_0=1
\]
in $(\widetilde\Pi_n^D)_{\mathrm{ab}}$
\[
\begin{aligned}
\tilde\sigma^{\unsure{\ast|-1}}
\bigl(\overset{\inserted{\unsure{\zeta}}}{\unsure{\tilde\rho}}_0\bigr)
&=\unsure{\zeta_1|\gamma_1}+\deleted{\unsure{\tilde\rho}}z_0,\\
\tilde\sigma^{\unsure{\ast|-1}}(\unsure{\zeta_1|\gamma_1})
&=\unsure{\zeta_0|\gamma_0}+z_1,\quad\cdots\\
\tilde\sigma^{\unsure{\ast|-1}}(\unsure{\zeta_i|\rho_i})
&=\unsure{\zeta_{n+1-i}|\gamma_{n+1-i}}+z_i
\end{aligned}
\]

$z_i=\tilde\omega_0^{\unsure{\nu_i|\gamma_i}}$
the condition for \unsure{invariance} of $\unsure{\widetilde\Pi_n^\circ|\Pi_n^\circ}$ is
that the $\unsure{\nu_i|\gamma_i}\equiv0\ (2)$.
We are going to \unsure{try} to \deleted{\unsure{them}} \unsure{calculate} at least $z_0,z_1$.

\margin{$\tilde\rho_0=\tilde\sigma\tilde\rho\tilde\sigma^{-1}\tilde\rho$}
\[
\tilde\sigma(\tilde\rho_0)
=\tilde\sigma^{\unsure{\ast|-1}}(\tilde\rho_0)
=\tilde\rho\tilde\sigma^{-1}\tilde\rho\tilde\sigma
\]
\[
\tilde\rho_1=\tilde\rho\tilde\sigma\tilde\rho\tilde\sigma^{-1}
\]
\[
\begin{aligned}
\tilde\rho_1^{-1}\tilde\sigma^{-1}(\tilde\rho_0)
&=\tilde\sigma\tilde\rho^{-1}
\underbrace{\underbrace{\tilde\sigma^{-1}
\underbrace{\tilde\rho^{-1}\tilde\rho}\tilde\sigma^{-1}}_{
\tilde\sigma^{-2}=\tilde\rho^{-n}}}_{\tilde\rho^{-n}}
\tilde\rho\tilde\sigma\\
&=\tilde\sigma\underbrace{\tilde\rho^{-n}}_{
\unsure{\tilde\omega_0^{-1}|\tilde\omega_1^{-1}}}\tilde\sigma
=\deleted{\unsure{z_0}}\tilde\omega_0^{-1}\underbrace{\tilde\sigma^2}_{\tilde\omega_0}=1
\end{aligned}
\]
\margin{\deleted{$\unsure{\tilde\rho}$}}

\deleted{$\tilde\sigma(\tilde\rho_1)$}
\[
\tilde\sigma^{\unsure{\ast|-1}}(\tilde\rho_1)
=\tilde\sigma^{-1}(\tilde\rho_1)
=\tilde\sigma^{-1}\tilde\rho\tilde\sigma\tilde\rho
\]
\[
\tilde\rho_0^{-1}\tilde\sigma(\tilde\rho_1)
=\unsure{\tilde\rho^{-1}}
\underbrace{\tilde\sigma^{\unsure{\ast|-1}}
\underbrace{\tilde\rho^{-1}
\underbrace{\unsure{\tilde\sigma^{-1}}\tilde\sigma^{-1}}_{
\omega_0^{-2}}\tilde\rho}_{\omega_0^{\unsure{-2|-1}}}
\tilde\sigma}_{1}\tilde\rho=1
\]

\deleted{$\unsure{\sigma|\rho}$}
\[
\tilde\rho_i=\tilde\rho^i\tilde\sigma\tilde\rho\tilde\sigma^{-1}\tilde\rho^{1-i}
\]
\[
\tilde\sigma(\tilde\rho_i)
=\tilde\sigma\tilde\rho^i\tilde\sigma\tilde\rho\tilde\sigma^{-1}\tilde\rho^{1-i}\tilde\sigma^{-1}
\]
\[
\tilde\rho_n\cdots\tilde\rho_{\unsure{n+2-i|n-2-i}}\tilde\rho_{n+1-i}\tilde\rho_{n+2-i}^{-1}
\]
\margin{5}